

\begin{theorem}
    Jeżeli wszystkie wyrazy ciągu $\boldsymbol{\left (a_{n} \right)}$ są równe tej 
    samej liczbie $\boldsymbol{m}$, tzn.: jeżeli ciąg $\boldsymbol{\left (a_{n} \right)}$
    jest ciągiem stałym, to liczba $\\$ $\boldsymbol{m}$ jest granicą tego ciągu.
\end{theorem}
\begin{proof}[\textbf{Dowód}]
    Ponieważ ciąg $\boldsymbol{\left (a_{n} \right)}$ jest ciągiem stałym to dla każdej 
    liczby naturalnej $\boldsymbol{n}$ jest $\boldsymbol{a_{n}=m}$. Z definicji mamy, że dla
    dowolnie małej liczby dodatniej $\boldsymbol{\epsilon}$ istnieje taka 
    liczba naturalna $\boldsymbol{N}$, że dla każdej liczby naturalnej $\boldsymbol{n>N}$
    zachodzi nierówność: $\boldsymbol{\left |a_{n}-m \right|<\epsilon}$.
    Ale ponieważ dla każdego $\boldsymbol{n}$ jest $\boldsymbol{a_{n}=m}$ to otrzymujemy
    $\boldsymbol{\left |a_{n}-m \right|=\left |m-m \right|=0<\epsilon}$. A zatem liczba 
    $\boldsymbol{m}$ jest granicą ciągu $\boldsymbol{\left (a_{n} \right)}$.
\end{proof}
